\chapter{Tree-Based Methods for Finance}
\label{ch:trees}

\begin{projectconnection}[Phase 2: Your Primary Model]
LightGBM is your primary model in Phase~2.
This chapter teaches the algorithm, the hyperparameter choices specific to financial data, and the Gu--Kelly--Xiu (2020) evidence that tree methods are competitive with---and often beat---deep learning for tabular financial data.
The key insight: aggressive regularization (low depth, high \texttt{min\_child\_samples}) is critical in finance's low signal-to-noise regime.
Every LightGBM result you report must appear alongside the ridge baseline from Chapter~\ref{ch:regularized}.
\end{projectconnection}

%% ═══════════════════════════════════════════════════════════
\section{Decision Trees}
\label{sec:decision-trees}
%% ═══════════════════════════════════════════════════════════

\begin{prereq}[Information Gain and Impurity Measures]
A \emph{splitting criterion} measures how ``mixed'' a node is.
For regression, the standard measure is variance (or equivalently, MSE).
For classification, the two workhorses are:
\begin{itemize}[nosep]
  \item \textbf{Gini impurity}: $G = \sum_{k=1}^K \hat{p}_k (1 - \hat{p}_k)$, where $\hat{p}_k$ is the fraction of class $k$ in the node. Gini is zero when all samples belong to one class (pure node) and maximized when classes are equally represented.
  \item \textbf{Entropy}: $H = -\sum_{k=1}^K \hat{p}_k \log_2 \hat{p}_k$. Same idea, different scale.
\end{itemize}
\textbf{Information gain} from a split is the parent impurity minus the weighted average of child impurities.
The algorithm picks the split that maximizes this gain.
In practice, Gini and entropy produce nearly identical trees; the choice rarely matters.
\end{prereq}

\subsection{CART: Recursive Binary Splitting}

\begin{definition}[Classification and Regression Tree (CART)]
A decision tree partitions the feature space $\R^p$ into $J$ non-overlapping rectangular regions $R_1, R_2, \ldots, R_J$ by a sequence of binary splits.
For regression, the prediction in region $R_j$ is the mean of the training responses in that region:
\begin{equation}
\hat{f}(\bx) = \sum_{j=1}^J c_j \cdot \mathbf{1}(\bx \in R_j), \qquad c_j = \frac{1}{|R_j|} \sum_{i: \bx_i \in R_j} y_i
\label{eq:cart-prediction}
\end{equation}
where $c_j$ is the leaf value (the average target in leaf $j$) and $\mathbf{1}(\bx \in R_j)$ is the indicator function.
\end{definition}

The tree is built top-down by \textbf{greedy recursive binary splitting}.
At each node, the algorithm considers every feature $j \in \{1, \ldots, p\}$ and every possible split point $s$, and selects the pair $(j, s)$ that minimizes the sum of squared residuals in the two child nodes:
\begin{equation}
\min_{j, s} \left[ \sum_{i: x_{ij} \leq s} (y_i - c_L)^2 + \sum_{i: x_{ij} > s} (y_i - c_R)^2 \right]
\label{eq:cart-split}
\end{equation}
where $c_L$ and $c_R$ are the means of the left and right child nodes.
This process repeats recursively until a stopping criterion is met (e.g., minimum samples per leaf, maximum depth, or minimum impurity decrease).

\subsection{Why Single Trees Overfit}

A fully grown tree (no stopping criterion) creates one leaf per training observation---it memorizes the training data perfectly.
The training MSE is zero, but the test MSE is terrible.
The fundamental problem: a single tree is a high-variance estimator.
Small perturbations in the training data produce completely different tree structures.

This is especially catastrophic in finance.
With monthly stock returns, the signal-to-noise ratio is roughly $0.05$ (the true signal explains $\sim$0.25\% of return variation; see the Gu--Kelly--Xiu numbers in Section~\ref{sec:gkx-horse-race}).
A single tree will partition on noise, find spurious splits, and produce worthless out-of-sample predictions.
No one uses single decision trees for return prediction.

%% ═══════════════════════════════════════════════════════════
\section{Random Forests}
\label{sec:random-forests}
%% ═══════════════════════════════════════════════════════════

\begin{prereq}[Bootstrap Aggregation (Bagging)]
Given a training set of $n$ observations, a \emph{bootstrap sample} draws $n$ observations \emph{with replacement}.
On average, each bootstrap sample contains $\sim$63.2\% of the unique training observations (the rest are duplicates).
\textbf{Bagging} fits a model to each of $B$ bootstrap samples and averages the predictions:
$\hat{f}_{\mathrm{bag}}(\bx) = \frac{1}{B}\sum_{b=1}^B \hat{f}_b(\bx)$.
Averaging reduces variance without increasing bias---the variance of the mean of $B$ i.i.d.\ estimates is $\sigma^2/B$.
The catch: if individual trees are correlated (because they all split on the same dominant feature), the variance reduction is limited.
\end{prereq}

Random forests (Breiman, 2001) combine bagging with a simple but powerful modification: at each split, only a random subset of $m$ features (out of $p$ total) is considered as split candidates.

\begin{keyidea}[Random Feature Subsets Decorrelate Trees]
If one feature is very strong, every bagged tree will split on it first, making the trees highly correlated.
Averaging correlated estimators provides limited variance reduction.
By restricting each split to a random $m \ll p$ features, the forest forces trees to explore different parts of the feature space, producing \emph{decorrelated} trees.
The variance of the forest average is:
\begin{equation}
\mathrm{Var}(\hat{f}_{\mathrm{RF}}) = \rho \sigma^2 + \frac{1-\rho}{B}\sigma^2
\label{eq:rf-variance}
\end{equation}
where $\rho$ is the average pairwise correlation between trees and $\sigma^2$ is the variance of a single tree.
As $B \to \infty$, the second term vanishes but the first does not---the irreducible variance floor is $\rho\sigma^2$.
Smaller $m$ reduces $\rho$ (more decorrelated trees) at the cost of slightly higher individual-tree bias.
The standard default is $m = \lfloor p/3 \rfloor$ for regression and $m = \lfloor \sqrt{p} \rfloor$ for classification.
\end{keyidea}

\subsection{Out-of-Bag Error}

Because each bootstrap sample omits $\sim$36.8\% of observations, each tree has a built-in holdout set.
The \textbf{out-of-bag (OOB) error} evaluates each observation using only the trees for which it was \emph{not} in the bootstrap sample:
\begin{equation}
\hat{f}_{\mathrm{OOB}}(\bx_i) = \frac{1}{|\mathcal{B}_i|} \sum_{b \in \mathcal{B}_i} \hat{f}_b(\bx_i)
\label{eq:oob-prediction}
\end{equation}
where $\mathcal{B}_i = \{b : \text{observation } i \text{ not in bootstrap sample } b\}$.
OOB error is approximately equivalent to leave-one-out cross-validation and comes for free---no separate validation split needed.

\begin{warning}[OOB Is Not Sufficient for Time Series]
OOB error assumes observations are exchangeable (i.i.d.).
In financial time series, this assumption fails: if observation $t{+}1$ is in a tree's training set and observation $t$ is OOB, the tree has seen the future.
For your project, use purged $K$-fold CV (Chapter~\ref{ch:validation}) instead of OOB for model selection.
OOB is still useful as a quick sanity check during development, but never as your primary validation metric.
\end{warning}

\subsection{Why Forests Are Better---but Still Limited}

Random forests reduce variance dramatically compared to single trees.
They are robust to hyperparameters (forests of 500+ trees with default $m$ generally work well) and handle mixed feature types, missing values, and interactions naturally.

But forests have two limitations for financial prediction:
\begin{enumerate}[nosep]
  \item \textbf{No sequential correction}: each tree is fit independently. If 100 trees all make the same systematic error, the 101st tree will too. There is no mechanism to focus on the ensemble's mistakes.
  \item \textbf{Equal weighting}: each tree contributes equally to the average, regardless of its accuracy.
\end{enumerate}
Gradient boosting addresses both limitations.

%% ═══════════════════════════════════════════════════════════
\section{Gradient Boosting}
\label{sec:gradient-boosting}
%% ═══════════════════════════════════════════════════════════

\begin{prereq}[Gradient Descent in Parameter Space]
Standard gradient descent updates parameters in the direction of the negative gradient: $\btheta_{m+1} = \btheta_m - \eta \nabla_{\btheta} \loss(\btheta_m)$.
Gradient boosting applies the same idea, but in \emph{function space}: instead of updating parameters, it adds a new function (a small tree) that points in the direction of the steepest descent of the loss.
\end{prereq}

\subsection{The Algorithm}

Gradient boosting builds an additive ensemble sequentially.
Starting from an initial constant prediction $F_0(\bx) = \bar{y}$, each iteration $m$ fits a small tree $h_m(\bx)$ to the \emph{negative gradient} of the loss function evaluated at the current ensemble's predictions:

\begin{equation}
F_m(\bx) = F_{m-1}(\bx) + \eta \cdot h_m(\bx)
\label{eq:gbm-update}
\end{equation}

Term by term:
\begin{itemize}[nosep]
  \item $F_{m-1}(\bx)$: the ensemble's prediction after $m{-}1$ trees. This is the running ``best guess.''
  \item $h_m(\bx)$: a shallow decision tree (typically depth 3--6) fit to the \emph{pseudo-residuals} $r_{im} = -\frac{\partial \loss(y_i, F_{m-1}(\bx_i))}{\partial F_{m-1}(\bx_i)}$. For squared-error loss, the pseudo-residuals are simply $r_{im} = y_i - F_{m-1}(\bx_i)$---the ordinary residuals.
  \item $\eta \in (0, 1]$: the \textbf{learning rate} (or shrinkage parameter). Smaller $\eta$ means each tree contributes less, requiring more trees but producing a smoother, more regularized ensemble. Typical values: $\eta \in [0.01, 0.1]$.
\end{itemize}

\begin{intuition}[Boosting as Iterative Refinement]
Think of boosting as iterative refinement.
The first tree captures the biggest, most obvious patterns in the data.
The second tree looks at what the first tree got wrong (the residuals) and corrects those errors.
The third tree corrects what the first two together still get wrong.
Each subsequent tree focuses on progressively smaller, harder-to-find patterns.

The learning rate $\eta$ controls how aggressively each correction is applied.
A small $\eta$ (say $0.01$) means each tree barely nudges the prediction---you need thousands of trees, but the ensemble generalizes better because no single tree dominates.
A large $\eta$ (say $1.0$) means each tree makes a full-strength correction---fast convergence but high risk of overfitting.
In finance, always err toward smaller $\eta$.
\end{intuition}

\subsection{Why Boosting Beats Bagging}

The key difference from random forests:
\begin{itemize}[nosep]
  \item \textbf{Forests}: each tree is fit independently on a random subsample. Trees don't communicate. The ensemble reduces variance by averaging.
  \item \textbf{Boosting}: each tree is fit \emph{sequentially} on the residuals of the current ensemble. Trees explicitly correct each other's errors. The ensemble reduces \emph{both} bias and variance.
\end{itemize}

This sequential error correction is why gradient boosting typically outperforms random forests for structured/tabular data.
Gu, Kelly, and Xiu (2020) found that gradient-boosted trees (GBRT) achieved a monthly OOS $R^2$ of $0.34\%$ on individual stock returns, substantially above random forests.

%% ═══════════════════════════════════════════════════════════
\section{XGBoost}
\label{sec:xgboost}
%% ═══════════════════════════════════════════════════════════

XGBoost (Chen and Guestrin, 2016) is gradient boosting with three key enhancements: a regularized objective, a second-order Taylor expansion for faster optimization, and systems-level engineering for speed.

\subsection{Regularized Objective}

XGBoost adds explicit regularization to the tree complexity.
The objective at iteration $m$ is:
\begin{equation}
\loss^{(m)} = \sum_{i=1}^n \ell\!\big(y_i,\; F_{m-1}(\bx_i) + h_m(\bx_i)\big) + \Omega(h_m)
\label{eq:xgb-objective}
\end{equation}
where the complexity penalty on tree $h_m$ is:
\begin{equation}
\Omega(h_m) = \gamma \, J + \frac{1}{2}\lambda \sum_{j=1}^J w_j^2
\label{eq:xgb-regularization}
\end{equation}

Term by term:
\begin{itemize}[nosep]
  \item $\ell(y_i, \hat{y}_i)$: the per-sample loss (e.g., squared error, log-loss).
  \item $\Omega(h_m)$: the regularization term penalizing tree complexity.
  \item $\gamma$: penalty per additional leaf. Larger $\gamma$ prunes the tree more aggressively. Acts like a minimum-gain threshold: a split must improve the loss by at least $\gamma$ to be accepted.
  \item $J$: the number of leaves in tree $h_m$.
  \item $\lambda$: $L_2$ penalty on the leaf weights $w_j$. Same idea as ridge regression (Chapter~\ref{ch:regularized}), but applied to leaf values rather than regression coefficients.
  \item $w_j$: the output value (prediction) in leaf $j$.
\end{itemize}

\subsection{Second-Order Approximation and Optimal Leaf Weights}

XGBoost applies a second-order Taylor expansion to the loss around the current prediction $F_{m-1}(\bx_i)$:
\begin{equation}
\loss^{(m)} \approx \sum_{i=1}^n \Big[ g_i \, h_m(\bx_i) + \frac{1}{2} h_i \, h_m(\bx_i)^2 \Big] + \Omega(h_m) + \text{const.}
\label{eq:xgb-taylor}
\end{equation}
where $g_i = \partial \ell / \partial F_{m-1}(\bx_i)$ is the first-order gradient and $h_i = \partial^2 \ell / \partial F_{m-1}(\bx_i)^2$ is the second-order gradient (Hessian) of the loss with respect to the current prediction.

For a given tree structure, the optimal leaf weight for leaf $j$ is:
\begin{equation}
w_j^* = -\frac{\sum_{i \in I_j} g_i}{\sum_{i \in I_j} h_i + \lambda}
\label{eq:xgb-leaf-weight}
\end{equation}
where $I_j$ is the set of training indices assigned to leaf $j$.
The denominator $\sum h_i + \lambda$ is the analog of $(\bX^\top\bX + \lambda\mathbf{I})$ from ridge---the Hessians measure curvature and $\lambda$ adds regularization.

\subsection{Split Gain}

XGBoost evaluates candidate splits using the gain formula:
\begin{equation}
\mathrm{Gain} = \frac{1}{2}\left[\frac{\big(\sum_{i \in I_L} g_i\big)^2}{\sum_{i \in I_L} h_i + \lambda} + \frac{\big(\sum_{i \in I_R} g_i\big)^2}{\sum_{i \in I_R} h_i + \lambda} - \frac{\big(\sum_{i \in I_P} g_i\big)^2}{\sum_{i \in I_P} h_i + \lambda}\right] - \gamma
\label{eq:xgb-gain}
\end{equation}

Term by term:
\begin{itemize}[nosep]
  \item The first two fractions: the objective value of the left and right children.
  \item The third fraction: the objective value of the parent (no split).
  \item $-\gamma$: the cost of adding one more leaf. If the gain from splitting is less than $\gamma$, the split is rejected---built-in pruning.
\end{itemize}

\subsection{Growth Strategy: Level-Wise}

XGBoost grows trees \textbf{level-wise} (breadth-first): it splits all nodes at the current depth before moving to the next depth level.
This produces balanced trees---every branch has the same depth.
Level-wise growth is safe (it doesn't chase isolated high-gain nodes) but can be wasteful: it may split nodes that contribute little, spending capacity on unimportant regions of the feature space.

%% ═══════════════════════════════════════════════════════════
\section{LightGBM}
\label{sec:lightgbm}
%% ═══════════════════════════════════════════════════════════

LightGBM (Ke et al., 2017) is an alternative gradient boosting implementation designed for speed and memory efficiency.
It introduces four key innovations over XGBoost: histogram-based splitting, leaf-wise growth, Gradient-based One-Side Sampling (GOSS), and Exclusive Feature Bundling (EFB).

\subsection{Histogram-Based Splitting}

Standard gradient boosting evaluates every unique feature value as a candidate split point.
With $n$ observations and $p$ features, this requires $O(n \cdot p)$ operations per level---expensive when $n$ is large.

LightGBM bins each continuous feature into $k$ discrete bins (default $k = 255$) using equal-frequency histograms.
Split finding then scans $k$ bins instead of $n$ unique values:
\begin{itemize}[nosep]
  \item \textbf{Complexity}: $O(k \cdot p)$ per level instead of $O(n \cdot p)$. Since $k = 255 \ll n$, this is much faster.
  \item \textbf{Memory}: storing bin indices (1 byte per feature per sample) instead of floating-point values (4--8 bytes) drastically reduces memory footprint.
  \item \textbf{Accuracy trade-off}: binning introduces a small discretization error, but Ke et al.\ (2017) show this has negligible effect on model quality. The regularization effect of binning can even help prevent overfitting.
\end{itemize}

\subsection{Leaf-Wise Growth}

\begin{keyidea}[Leaf-Wise vs.\ Level-Wise Growth]
XGBoost grows trees \textbf{level-wise}: split all nodes at the current depth, then proceed to the next depth.
This produces balanced trees but wastes splits on low-gain nodes.

LightGBM grows trees \textbf{leaf-wise}: at each step, find the leaf with the highest gain \emph{across the entire tree} and split it.
This concentrates the tree's limited capacity on the most informative regions of the feature space.

Leaf-wise growth typically produces lower training loss with the same number of leaves.
The risk: it can create very deep, unbalanced trees that overfit.
This is why \texttt{max\_depth} and \texttt{num\_leaves} are critical hyperparameters for LightGBM---you must constrain depth explicitly.
\end{keyidea}

\subsection{GOSS: Gradient-Based One-Side Sampling}

Gradient-based One-Side Sampling exploits a key observation: samples with large gradients (i.e., large residuals) are the ones the model gets most wrong and therefore carry the most information about which splits will reduce loss.

The GOSS algorithm:
\begin{enumerate}[nosep]
  \item Sort all samples by the absolute value of their gradient $|g_i|$.
  \item Keep the top $a \times 100\%$ of high-gradient samples (default: $a = 0.2$, keep 20\%).
  \item Randomly sample $b \times 100\%$ of the remaining low-gradient samples (default: $b = 0.1$).
  \item Up-weight the low-gradient subsample by a factor of $(1 - a)/b$ to maintain the correct gradient distribution.
\end{enumerate}

This reduces the effective sample size for split finding to $(a + b) \times n$ while preserving the gradient distribution.
Ke et al.\ (2017) show that GOSS achieves almost the same split quality as full-data scanning with a fraction of the computation.

\subsection{EFB: Exclusive Feature Bundling}

Many real-world feature sets contain mutually exclusive features---features that rarely take nonzero values simultaneously (e.g., one-hot encoded categoricals, sparse indicator features).
EFB identifies such features and bundles them into a single composite feature, reducing the effective number of features.

The algorithm:
\begin{enumerate}[nosep]
  \item Build a conflict graph: two features conflict if they frequently take nonzero values for the same sample.
  \item Greedily bundle features with few conflicts (graph coloring heuristic).
  \item Merge bundled features by offsetting their bin ranges so no information is lost.
\end{enumerate}

Finding the optimal bundling is NP-hard, but the greedy approximation works well in practice.
EFB is most impactful on high-dimensional sparse datasets.
For your project's 10--20 risk-cube features, EFB matters less---but it is part of what makes LightGBM fast by default.

\subsection{Speed: LightGBM vs.\ XGBoost}

Ke et al.\ (2017) report that LightGBM achieves up to $20\times$ speedup over XGBoost on large datasets while reaching the same accuracy.
The speedup comes from the combination of histogram binning ($O(k)$ vs.\ $O(n)$ per feature), leaf-wise growth (fewer wasted splits), and GOSS (fewer samples scanned).

For your project, the speed difference is less dramatic ($\sim$1,250 daily observations is tiny by ML standards), but LightGBM remains the preferred choice because:
\begin{enumerate}[nosep]
  \item The Python API (\texttt{lightgbm}) is mature and well-documented.
  \item Built-in handling of categorical features (no need for one-hot encoding).
  \item Native integration with SHAP for feature importance (Chapter~\ref{ch:interpretation}).
  \item It is the industry default for tabular data at quant desks and Kaggle competitions alike.
\end{enumerate}

%% ═══════════════════════════════════════════════════════════
\section{LightGBM vs.\ XGBoost: Comparison}
\label{sec:lgb-vs-xgb}
%% ═══════════════════════════════════════════════════════════

\begin{center}
\begin{tabular}{p{4.2cm} p{5cm} p{5cm}}
\toprule
\textbf{Feature} & \textbf{XGBoost} & \textbf{LightGBM} \\
\midrule
Growth strategy & Level-wise (balanced) & Leaf-wise (greedy, deeper) \\[4pt]
Split finding & Exact or approximate (quantile sketch) & Histogram-based ($k{=}255$ bins) \\[4pt]
Data sampling & Row subsampling & GOSS (gradient-based) \\[4pt]
Feature reduction & Column subsampling & Column subsampling + EFB \\[4pt]
Regularization & $L_1$ + $L_2$ on leaf weights, $\gamma$ pruning & Same, plus \texttt{min\_child\_samples} \\[4pt]
Categorical features & Requires encoding & Native handling \\[4pt]
Speed (large $n$) & Slower & Up to $20\times$ faster \\[4pt]
Overfitting risk & Lower (balanced trees) & Higher (must tune \texttt{max\_depth}) \\[4pt]
Default choice & Mature, well-tested & Industry standard for tabular data \\
\bottomrule
\end{tabular}
\end{center}

\medskip
\noindent\textbf{Bottom line}: For your project, use LightGBM as the primary model.
XGBoost is a fine alternative if you encounter LightGBM-specific issues (e.g., package installation problems on the GS compute environment).
The two produce similar accuracy when tuned well; the choice is primarily about speed and API convenience.

%% ═══════════════════════════════════════════════════════════
\section{Hyperparameters for Financial Data}
\label{sec:hyperparams-finance}
%% ═══════════════════════════════════════════════════════════

\begin{warning}[Never Use Default LightGBM Parameters on Financial Data]
LightGBM's defaults are tuned for Kaggle-style competitions with high signal-to-noise ratios and $10^5$--$10^7$ samples.
Financial return prediction has SNR $\sim 0.05$ and $\sim 10^3$ samples.
Default parameters (\texttt{num\_leaves=31}, \texttt{max\_depth=-1}, \texttt{min\_child\_samples=20}) will overfit catastrophically: you will see excellent in-sample $R^2$ and negative out-of-sample $R^2$.
Always start from the conservative grid below and tune from there.
\end{warning}

The table below gives recommended starting ranges for financial return prediction:

\begin{center}
\begin{tabular}{p{4.5cm} p{2.5cm} p{2cm} p{5cm}}
\toprule
\textbf{Parameter} & \textbf{Finance range} & \textbf{Default} & \textbf{Why} \\
\midrule
\texttt{max\_depth} & 3--5 & $-1$ (no limit) & Shallow trees = less overfitting. Depth 4 allows up to 3-way interactions. \\[4pt]
\texttt{num\_leaves} & 8--31 & 31 & Should be $\leq 2^{\texttt{max\_depth}}$. \\[4pt]
\texttt{min\_child\_samples} & 50--200 & 20 & Forces each leaf to contain enough samples for statistical reliability. With $n = 1{,}250$, a leaf of 20 is $1.6\%$ of data---noise. \\[4pt]
\texttt{learning\_rate} & 0.01--0.05 & 0.1 & Slow learning + early stopping. Lower $\eta$ with more trees. \\[4pt]
\texttt{n\_estimators} & 500--5000 & 100 & Paired with low \texttt{learning\_rate}. Use early stopping on validation loss. \\[4pt]
\texttt{subsample} & 0.6--0.8 & 1.0 & Row subsampling (per tree). Reduces variance. \\[4pt]
\texttt{colsample\_bytree} & 0.6--0.8 & 1.0 & Feature subsampling (per tree). Decorrelates trees, like random forests. \\[4pt]
\texttt{reg\_alpha} ($L_1$) & 0--0.1 & 0 & Mild $L_1$ on leaf weights. Usually less important than depth/samples. \\[4pt]
\texttt{reg\_lambda} ($L_2$) & 1--10 & 0 & $L_2$ on leaf weights. Shrinks leaf predictions toward zero. \\
\bottomrule
\end{tabular}
\end{center}

\begin{keyidea}[The Core Principle: Constrain Capacity, Not Just Trees]
In finance, regularization must be multi-layered:
\begin{enumerate}[nosep]
  \item \textbf{Small trees}: \texttt{max\_depth}~3--5 limits the complexity of each tree.
  \item \textbf{Slow learning}: \texttt{learning\_rate}~0.01--0.05 means each tree barely nudges the ensemble.
  \item \textbf{Subsampling}: \texttt{subsample} and \texttt{colsample\_bytree}~0.6--0.8 inject randomness.
  \item \textbf{Leaf regularization}: \texttt{min\_child\_samples}~50--200 prevents tiny, noisy leaves.
  \item \textbf{Early stopping}: monitor validation loss, stop when it plateaus or increases.
\end{enumerate}
No single knob is sufficient.
You need all five simultaneously because the signal is so weak that overfitting can occur through any unguarded channel.
\end{keyidea}

\begin{workedexample}{Tuning LightGBM for VaR Utilization Signals}
Suppose you have $n = 1{,}250$ daily observations and $p = 15$ risk-cube features.
You want to predict next-day VIX innovations using VaR utilization and factor concentration features.

\textbf{Step 1: Set a conservative base.}
\begin{verbatim}
params = {
    'max_depth': 4,
    'num_leaves': 15,
    'min_child_samples': 100,
    'learning_rate': 0.02,
    'n_estimators': 3000,
    'subsample': 0.7,
    'colsample_bytree': 0.7,
    'reg_lambda': 5.0,
    'early_stopping_rounds': 50,
}
\end{verbatim}

\textbf{Step 2: Early stopping.}
Use the purged validation fold (Chapter~\ref{ch:validation}) as the early-stopping monitor.
Training will stop when validation MSE has not improved for 50 rounds.
With \texttt{learning\_rate=0.02}, you might stop at $\sim$800 trees out of the 3,000 maximum.

\textbf{Step 3: Compare to ridge.}
If this LightGBM's purged-CV $R^2$ does not exceed the ridge baseline from Chapter~\ref{ch:regularized} on the same features, reduce \texttt{max\_depth} to 3 or increase \texttt{min\_child\_samples} to 150.
If it still doesn't beat ridge, the signal may be approximately linear---and that is a valid finding.
\end{workedexample}

%% ═══════════════════════════════════════════════════════════
\section{The GKX Horse Race}
\label{sec:gkx-horse-race}
%% ═══════════════════════════════════════════════════════════

Gu, Kelly, and Xiu (2020, RFS) ran the most comprehensive ML horse race in empirical asset pricing: 30,000 individual US stocks, 60 years (1957--2016), 94 stock characteristics plus 8 macro variables plus 74 industry dummies ($\sim$900 predictors after interactions), with monthly return prediction evaluated on 30 years of out-of-sample data.

\begin{keyresult}[Gu--Kelly--Xiu (2020): The ML Horse Race]
Monthly out-of-sample $R^2$ for individual stock return prediction (value-weighted):
\begin{center}
\begin{tabular}{lc}
\toprule
\textbf{Model} & \textbf{Monthly OOS $R^2$} \\
\midrule
OLS (all 900+ predictors) & $< 0$ (negative) \\
Elastic net (ENet+H) & $0.11\%$ \\
PLS & $0.27\%$ \\
PCR & $0.26\%$ \\
Random forest & $0.28\%$ \\
Gradient-boosted trees (GBRT) & $\mathbf{0.34\%}$ \\
Neural network (NN3, 3 hidden layers) & $\mathbf{0.40\%}$ \\
\bottomrule
\end{tabular}
\end{center}

Portfolio-level results: the NN3 long-short decile portfolio earned an annualized Sharpe ratio of $1.35$ (value-weighted).
Trees and neural nets achieve their edge through nonlinear predictor interactions that linear methods miss.
\end{keyresult}

\subsection{What the Numbers Mean}

\begin{enumerate}[nosep]
  \item \textbf{OLS collapses}: with 900+ predictors and $\sim$30,000 stocks/month, OLS massively overfits. The OOS $R^2$ is \emph{negative}---worse than predicting the mean. This is the curse of dimensionality in action.

  \item \textbf{Regularization rescues linear models}: elastic net shrinks the 900+ coefficients, pulling OOS $R^2$ to $0.11\%$. PLS and PCR do even better ($0.26$--$0.27\%$) by projecting onto a few linear combinations.

  \item \textbf{Trees beat linear methods}: GBRT at $0.34\%$ is $3\times$ elastic net's $0.11\%$. The gain comes from capturing nonlinear interactions between characteristics that linear methods miss.

  \item \textbf{Deep learning offers a modest further gain}: NN3's $0.40\%$ is ${\sim}18\%$ above GBRT's $0.34\%$. The margin is real but not enormous---and neural networks are harder to tune, interpret, and deploy.

  \item \textbf{Tiny $R^2$, real money}: $0.40\%$ monthly $R^2$ means the model explains less than half a percent of return variation. Yet it produces a long-short Sharpe ratio of $1.35$---excellent by asset management standards. In finance, tiny predictability translates to economically meaningful portfolio performance.
\end{enumerate}

\begin{projectconnection}[Why Trees Are Your Default]
For your project, GBRT (via LightGBM) is the right default model for three reasons:
\begin{enumerate}[nosep]
  \item \textbf{Competitive accuracy}: GBRT's $0.34\%$ monthly $R^2$ is close to the best neural network ($0.40\%$). On your smaller, lower-dimensional feature set (10--20 risk-cube features vs.\ GKX's 900+), the gap will likely shrink further.
  \item \textbf{Interpretability}: LightGBM integrates natively with SHAP (Chapter~\ref{ch:interpretation}), giving you feature importance that you can present to the desk. Neural networks require more work to interpret.
  \item \textbf{Robustness}: trees handle mixed feature types, missing data, and feature interactions without the extensive preprocessing that neural networks require. With $\sim$1,250 observations, you cannot afford the hyperparameter search that deep learning demands.
\end{enumerate}
If LightGBM does not beat ridge, the signal is linear (a valid finding).
Only if both ridge and LightGBM produce IC $> 0$ should you consider whether neural networks add value---and even then, the GKX evidence suggests the incremental gain is modest.
\end{projectconnection}

%% ═══════════════════════════════════════════════════════════
\section{When Trees Fail}
\label{sec:trees-fail}
%% ═══════════════════════════════════════════════════════════

Trees are not a universal solution.
They fail in specific, predictable ways that matter for financial applications.

\subsection{Extrapolation}

Trees cannot extrapolate beyond the range of training data.
A tree's prediction is always a weighted average of training-set target values---it can only output values it has seen.
If VaR utilization hits 98\% but the training data never exceeded 85\%, the tree will predict the same value as for 85\%.

This matters in finance because extreme events are precisely when you most need accurate predictions.
Linear models (including ridge) can extrapolate along the learned coefficient direction, which is at least directionally correct.
Trees are capped at the training-set boundary.

\textbf{Mitigation}: include recent extreme events in the training window when possible, and supplement tree predictions with linear models for tail scenarios.

\subsection{Non-Stationarity and Regime Changes}

Trees learn fixed split thresholds.
If the split ``VaR utilization $> 75\%$'' was informative in 2019--2023, the tree will apply it forever.
But if the risk limit structure changes in 2024 (e.g., the desk is restructured), 75\% may no longer be the relevant threshold.

\textbf{Mitigation}: rolling-window retraining with an expanding or sliding window (see backtesting in Chapter~\ref{ch:backtesting}).
Re-estimate the model every month or quarter using only recent data.
Regime conditioning (Chapter~\ref{ch:regimes}) can further help by allowing the model to adapt to changing market environments.

\subsection{Small Samples}

With $n = 1{,}250$ daily observations and depth-4 trees, the deepest leaves partition the data into at most $2^4 = 16$ bins.
If \texttt{min\_child\_samples=100}, you have at most $1{,}250 / 100 = 12$ effective leaves---the tree is heavily constrained.
This is by design (overfitting prevention), but it limits the complexity of interactions the tree can capture.

\textbf{Mitigation}: panel structure (Chapter~\ref{ch:panel}).
If you pool multiple asset classes, your effective sample size grows.
With 5 asset classes $\times$ 1,250 days $= 6{,}250$ observations, you can afford slightly deeper trees or lower \texttt{min\_child\_samples}---but always validate with purged CV.

\subsection{Comparison: Where Trees vs.\ Linear Models Win}

\begin{center}
\begin{tabular}{p{4cm} p{5cm} p{5cm}}
\toprule
\textbf{Scenario} & \textbf{Trees win} & \textbf{Linear models win} \\
\midrule
Feature interactions & Trees capture them automatically & Must engineer interactions by hand \\[4pt]
Nonlinear effects & Natural (via recursive splits) & Must add polynomial terms \\[4pt]
Extrapolation & Fail (bounded by training range) & Extrapolate along coefficients \\[4pt]
Small $n$, low SNR & Overfit risk; need heavy tuning & More stable (fewer parameters) \\[4pt]
Interpretability & SHAP, but complex & Coefficients are directly readable \\[4pt]
Missing data & Handle natively & Require imputation \\
\bottomrule
\end{tabular}
\end{center}

%% ═══════════════════════════════════════════════════════════
\section{Summary}
\label{sec:trees-summary}
%% ═══════════════════════════════════════════════════════════

\begin{center}
\begin{tabular}{p{4cm} p{9.5cm}}
\toprule
\textbf{Concept} & \textbf{Key Takeaway} \\
\midrule
Decision trees (CART) & Recursive binary splitting on features. Single trees overfit catastrophically in low-SNR regimes. Never use alone. \\[6pt]
Random forests & Bagging + random feature subsets decorrelate trees, reducing variance. OOB error is free but invalid for time series. \\[6pt]
Gradient boosting & Sequential: each tree corrects the ensemble's residuals. Reduces both bias and variance. Learning rate $\eta$ controls regularization. \\[6pt]
XGBoost & Level-wise growth, built-in $L_1$/$L_2$ on leaf weights, second-order Taylor expansion for split gain. \\[6pt]
LightGBM & Leaf-wise growth, histogram binning ($k{=}255$), GOSS (gradient-based sampling), EFB (feature bundling). Up to $20\times$ faster than XGBoost. \\[6pt]
Financial hyperparameters & \texttt{max\_depth}~3--5, \texttt{min\_child\_samples}~50--200, \texttt{subsample}~0.6--0.8, \texttt{learning\_rate}~0.01--0.05. \textbf{Never use defaults.} \\[6pt]
GKX horse race & GBRT $R^2 = 0.34\%$, best NN $R^2 = 0.40\%$ monthly OOS (Gu--Kelly--Xiu, 2020). Trees are the right default for tabular financial data. \\[6pt]
When trees fail & Cannot extrapolate, brittle to regime changes, need rolling retraining. Supplement with linear models for tail scenarios. \\
\bottomrule
\end{tabular}
\end{center}

\bigskip
\noindent\textbf{What's next}: Chapter~\ref{ch:interpretation} covers model interpretation---SHAP values, feature importance stability (MDA across folds), and how to present LightGBM results to a desk that thinks in terms of economic intuition, not tree splits.
